%------------------------------------------------------------------------------%
\documentclass[fleqn,utf8,aspectratio=169,12pt,ignorenonframetext]{beamer}

\usepackage{calculo_varias_variaveis-1}

\title{Números Complexos}

%------------------------------------------------------------------------------%
\begin{document}

\addtitle

%------------------------------------------------------------------------------%
\section{Unidade Imaginária}

%------------------------------------------------------------------------------%
\begin{frame}{Motivação}
  Considere a equação
  \[ x^2 + 1 = 0 \]
  \pause
  Manipulando formalmente temos a solução
  \[ x = \sqrt{-1} \]
  \pause
  Porém, \(x^2 \geq 0\) para todo $x\in\R$

  \pause
  \vspace{\baselineskip}
  Não existe \(\sqrt{-1}\) dentro dos números reais
\end{frame}

%------------------------------------------------------------------------------%
\begin{frame}{Unidade imaginária}
  \emph{Unidade imaginária}
  \[
    i^2 = -1
    \qquad\text{ou}\qquad
    i = \sqrt{-1}
  \]

  \pause
  A Engenharia Elétrica usa
  \vspace{\baselineskip}
  \(j^2=-1\)
\end{frame}

%------------------------------------------------------------------------------%
\begin{frame}{Propriedades de Potências}
  \large
  \[
    a^p a^q = a^{p+q}
  \]

  \pause
  \vspace{0.\baselineskip}
  \[
    \frac{a^p}{a^q} = a^{p-q}
  \]

  \pause
  \vspace{0.\baselineskip}
  \[
    \left(a^p\right)^q = a^{pq}
  \]
\end{frame}

%------------------------------------------------------------------------------%
\begin{frame}{Potências de $i$}
\begin{columns}
\column{0.5\linewidth}
  \begin{align*}
    \onslide<+->{i^1 & = i \\}
    \onslide<+->{i^2 & = -1 \\}
    \onslide<+->{i^3}
    \onslide<+->{& = i^2 \, i \;}
    \onslide<+->{ = (-1) i}
    \onslide<+->{ = -i}   \\
    \onslide<+->{i^4}
    \onslide<+->{ & = i^2 \, i^2}
    \onslide<+->{  = (-1)(-1)}
    \onslide<+->{ = 1}  \\
    \onslide<+->{i^5}
    \onslide<+->{ & = i^4 \, i \;}
    \onslide<+->{ = 1 i}
    \onslide<+->{ = i} \\
    \onslide<+->{i^6}
    \onslide<+->{ & = i^4 \, i^2}
    \onslide<+->{ = 1(-1)}
    \onslide<+->{ = -1}    \\
    \onslide<+->{i^7}
    \onslide<+->{ & = i^4 \, i^3}
    \onslide<+->{ = 1(-i)}
    \onslide<+->{ = -i}    \\
    \onslide<+->{i^8}
    \onslide<+->{ & = i^4 \, i^4}
    \onslide<+->{  = 1\times 1}
    \onslide<+->{ = 1}
  \end{align*}
\column{0.5\linewidth}
\begin{align*}
  \onslide<+->{i^0 & = 1\\}
  \onslide<+->{i^{-1}}
  \onslide<+->{& = \frac{1}{i}}
  \onslide<+->{  = \frac{i}{i^2}}
  \onslide<+->{  = \frac{i}{-1}}
  \onslide<+->{  = -i\\}
  \onslide<+->{i^{-2}}
  \onslide<+->{& = \frac{1}{i^2}}
  \onslide<+->{  = \frac{1}{-1}}
  \onslide<+->{  = -1\\}
  \onslide<+->{i^{-3}}
  \onslide<+->{& = \frac{1}{i^3}}
  \onslide<+->{  = \frac{1}{-i}}
  \onslide<+->{  = \frac{i}{-i^2}}
  \onslide<+->{  = \frac{i}{1}}
  \onslide<+->{  = i\\}
\end{align*}

\[
\onslide<+->{i^{75} =}
\onslide<+->{?}
\]
\end{columns}
\end{frame}

%------------------------------------------------------------------------------%
\begin{frame}{Potências de $i$}
\begin{columns}
\column{0.5\linewidth}
  \onslide<+->{Para uma potência $n$ inteira}

  \onslide<+->{
  \vspace{\baselineskip}
  Dividimos $n$ por 4}
  \[
    \onslide<+->{
    \begin{array}{>{\;}c<{\;}>{\;}c<{\;}}
      n & \multicolumn{1}{|c}{4} \\ \cline{2-2}
      r & q
    \end{array}}
    \qquad
    \onslide<+->{
    \begin{array}{>{\;}c<{\;}>{\;}c<{\;}}
      75 & \multicolumn{1}{|c}{4} \\ \cline{2-2}
      3  & 18
    \end{array}}
  \]
  \onslide<+->{Assim}
  \begin{gather*}
    \onslide<+->{n = 4q + r} \\
    \onslide<+->{75 = 4\times 18 + 3}
  \end{gather*}
\column{0.5\linewidth}
{\large
  \begin{align*}
    \onslide<+->{i^n}
    \onslide<+->{& = i^{\,4q\,+\,r} \\}
    \onslide<+->{& = i^{\,4q} \; i^r \\}
    \onslide<+->{& = \left(i^4\right)^q \; i^r \\}
    \onslide<+->{& = \left(1\right)^q \; i^r \\}
    \onslide<+->{& = i^r}
  \end{align*}}
  \onslide<+->{
  \[
    i^{75} = i^3 = -i
  \]}
\end{columns}
\end{frame}

%------------------------------------------------------------------------------%
\begin{frame}{Cuidado com raízes}
\begin{columns}
\column{0.4\linewidth}
  Sabemos que
  \[
    \sqrt{a}\,\sqrt{b} = \sqrt{ab}
  \]

  \pause
  Por exemplo
  \[
    \sqrt{9}\,\sqrt{4}
    \pause
    = 3\times 2
    \pause
    = 6
  \]
  \[
    \pause
    \sqrt{9\times 4}
    \pause
    = \sqrt{36}
    \pause
    = 6
  \]
\column{0.6\linewidth}
  \pause
  \[
    \sqrt{-9}\,\sqrt{-4}
    \pause
    = 3i\times 2i
    \pause
    = 6i^2
    \pause
    = {\color<15->{red}-6}
  \]
  \[
    \pause
    \sqrt{(-9)(-4)}
    \pause
    = \sqrt{36}
    \pause
    = {\color<15->{red}6}
  \]

  \vspace{\baselineskip}
  \onslide<16>{\emph{Só é verdade para $a$ e $b$ positivos!}}
\end{columns}
\end{frame}

%------------------------------------------------------------------------------%
%------------------------------------------------------------------------------%
\begin{question}
  Simplifique as expressões
  \Large
  \[
    a = i^7
  \]

  \[
    b = i^{57}
  \]

  \[
    c = \frac{i^5 + i^{17}}{i^3}
  \]
\end{question}

%------------------------------------------------------------------------------%
\begin{resolution}{Solução}
  \Large
  \[
    a
    \;=\; i^7
    \pause
    \;=\; i^4 \, i^3
    \pause
    \;=\; i^3
    \pause
    \;=\; i^2 i
    \pause
    \;=\; -i
  \]
\end{resolution}

%------------------------------------------------------------------------------%
\begin{resolution}{Solução}
  \Large
  \[
    b
    \;=\; i^{57}
    \pause
    \;=\; i^{4\times 14} \, i^1
    \pause
    \;=\; i
  \]
\end{resolution}

%------------------------------------------------------------------------------%
\begin{resolution}{Solução}
  \Large
  \[
    c
    \;=\; \frac{i^5 + i^{17}}{i^3}
    \pause
    \;=\; \frac{i^5}{i^3} + \frac{i^{17}}{i^3}
    \pause
    \;=\; i^{5-3} + i^{17-3}
  \]

  \pause
  \[
    \vphantom{c}
    \;=\; i^2 + i^{14}
    \pause
    \;=\; -1 + i^{\,4\times 3+2}
    \pause
    \;=\; -1 + i^2
  \]

  \pause
  \[
    \vphantom{c}
    \;=\; -1 -1
    \pause
    \;=\; -2
  \]
\end{resolution}

%------------------------------------------------------------------------------%


%------------------------------------------------------------------------------%
\section{Números Complexos}

%------------------------------------------------------------------------------%
\begin{frame}{Números complexos}
  Conjunto dos Números Complexos
  \[
    \mathbb{C} = \big\{\, z=a+bi \;\big|\; a, b \in \R \,\big\}
  \]
  \pause
  Parte real de $z$
  \[
    \Re(z) = a
  \]
  \pause
  Parte imaginária de $z$
  \[
    \Im(z) = b
  \]
  \pause
  Os números complexos com parte imaginária nula são os Reais
\end{frame}

%------------------------------------------------------------------------------%
\begin{frame}{Igualdade entre números complexos}

  \(z_1=a+bi\) \; é \emph{igual} a \; \(z_2=c+di\)

  \pause
  \vspace{\baselineskip}
  se, e somente se,
  \[
     a = c
     \qquad\text{e}\qquad
     b = d
  \]
\end{frame}

%------------------------------------------------------------------------------%
\begin{frame}{Soma e produto de números complexos}
  \onslide<+->{\emph{Soma e subtração}}
  \[
    \onslide<+->{(a+b\emph{i}) + (c+d\emph{i})}
    \onslide<+->{\;=\;  a + b\emph{i} + c + d\emph{i}}
    \onslide<+->{\;=\; (a+c) + (b+d) \emph{i}}
  \]
  \vspace{-\baselineskip}
  \[
    \onslide<+->{(a+b\emph{i}) - (c+d\emph{i})}
    \onslide<+->{\;=\; a+b\emph{i} -c -d\emph{i}}
    \onslide<+->{\;=\; (a-c) + (b-d) \emph{i}}
  \]

  \onslide<+->{\emph{Produto}}
  \begin{align*}
    \onslide<+->{
      (a+b\emph{i}) \times (c+d\emph{i})}
    \onslide<+->{
      & \;=\; a         \times (c+d\emph{i})
        \;+\; b\emph{i} \times (c+d\emph{i}) } \\
    \onslide<+->{
      & \;=\; ac + ad\emph{i} + bc\emph{i} + bd\emph{i}^2 } \\
    \onslide<+->{
      & \;=\; ac + ad\emph{i} + bc\emph{i} - bd } \\
    \onslide<+->{
      & \;=\; (ac-bd) \;+\; (ad+bc) \emph{i} }
  \end{align*}
\end{frame}

%------------------------------------------------------------------------------%
\begin{frame}{Conjugado}
  O \emph{conjugado} do número complexo
  \[
    z = a + bi
  \]
  \pause
  é o número
  \[
    \bar{z} = a - bi
  \]
\end{frame}

%------------------------------------------------------------------------------%
\begin{frame}{Propriedades dos conjugados}
  \begin{enumerate}[<+->]
  \setlength\itemsep{1.5em}

  \item \;
  \(
    \bar{\bar{z}} \onslide<+->{= z}
  \)

  \item \; se \(\bar{z}=z\) \onslide<+->{então $z$ é real}

  \item \;
  \(
    z\bar{z}
    \onslide<+->{= (a+bi)(a-bi)}
    \onslide<+->{= a^2 -abi +abi - b^2 i^2}
    \onslide<+->{= a^2 + b^2}
  \)
  \end{enumerate}
\end{frame}

%------------------------------------------------------------------------------%
\begin{frame}{Divisão de números complexos}
  \onslide<+->{\emph{Divisão}}%
  \onslide<+->{: um número $q$ é a divisão de \(z_1=a+bi\) por \(z_2=c+di\) se}
  \begin{columns}
  \column{0.5\baselineskip}
    \begin{align*}
      \onslide<+->{q z_2           & = z_1                   \\}
      \onslide<+->{q z_2 \bar{z}_2 & = z_1 \bar{z}_2         \\}
      \onslide<+->{q (c^2 + d^2)   & = (c+di)(a-bi)          \\}
      \onslide<+->{                & = ca -cbi +dai -dbi^2   \\}
      \onslide<+->{                & = ca -cbi +dai +dbi     \\}
      \onslide<+->{                & = (ac + bd) + (ad -bc)i}
    \end{align*}
  \column{0.5\baselineskip}
    \begin{align*}
      \onslide<+->{q & = \frac{(ac + bd) + (ad -bc)i}{c^2 + d^2}           \\[2mm]}
      \onslide<+->{  & = \frac{(ac + bd)}{c^2 + d^2} + \frac{(ad -bc)}{c^2 + d^2}i}
    \end{align*}
  \end{columns}
\end{frame}

%------------------------------------------------------------------------------%
\begin{frame}{Operações com os Números Complexos}
  Divisão de \(z_1=a+bi\) por \(z_2=c+di\)
  \[
    \pause
    \frac{z_1}{z_2}
    \pause
    \;=\; \frac{z_1\,\bar{z}_2}{z_2\,\bar{z}_2}
    \pause
    \;=\; \frac{(ac + bd)}{c^2 + d^2} + \frac{(ad -bc)}{c^2 + d^2}i
  \]
\end{frame}

%------------------------------------------------------------------------------%
\begin{frame}{Corpo complexo}
  A soma e o produto nos complexos tem as mesmas propriedades dos reais

  \vspace{\baselineskip}
  Por isso o conjunto dos complexos é chamado \emph{corpo}

  \vspace{\baselineskip}
  Exemplos de corpos: racionais, reais, complexos
\end{frame}

%------------------------------------------------------------------------------%
\include{L-01-Numeros_Complexos-ex2}

%------------------------------------------------------------------------------%
\section{Plano Complexo}

%------------------------------------------------------------------------------%
\begin{frame}{Plano Complexo}
\begin{columns}[b]
\column{0.55\linewidth}
  \begin{gather*}
    \onslide<1->{z = x + yi \quad          \in \mathbb{C} \\}
    \onslide<1->{(x, y)     \hspace{7.2ex} \in \R^2}
  \end{gather*}
  \onslide<2->{Parte real \tabto{24mm} \( \Re(z) = x \)       }\\[4mm]
  \onslide<3->{Parte imag \tabto{24mm} \( \Im(z) = y \)       }\\[4mm]
  \onslide<4->{Módulo     \tabto{24mm}}%
  \(
    \onslide<4->{\rho}
    \onslide<5->{= \abs{z}}
    \onslide<6->{= \sqrt{z\bar{z}}}
    \onslide<7->{= \sqrt{x^2+y^2\,}}
  \) \\[4mm]
  \onslide<8->{Argumento \tabto{24mm}}%
  \(
    \onslide<8->{\varphi}
    \onslide<9->{= \arg{(z)}}
    \onslide<10->{= \arctan\left(\frac{\,y\,}{x}\right)}
  \) \\[4mm]
  \onslide<11->{Conjugado  \tabto{24mm} \( \bar{z} = x - yi \) }
\column{0.3\linewidth}
\begin{center}%
  \only<1|handout:0>   {\includegraphics{plano_complexo-1}}% base
  \only<2|handout:0>   {\includegraphics{plano_complexo-2}}% real
  \only<3|handout:0>   {\includegraphics{plano_complexo-3}}% imaginario
  \only<4-7|handout:0> {\includegraphics{plano_complexo-4}}% modulo
  \only<8-10|handout:0>{\includegraphics{plano_complexo-5}}% argumento
  \only<11|handout:0>  {\includegraphics{plano_complexo-6}}% conjugado
  \only<12>            {\includegraphics{plano_complexo-7}}% conjugado
\end{center}
\end{columns}
\end{frame}

%------------------------------------------------------------------------------%
\include{L-01-Numeros_Complexos-ex3}

%------------------------------------------------------------------------------%
\section{Lista Mínima}

%------------------------------------------------------------------------------%
\begin{frame}{Lista Mínima}
  Atenção: A prova é baseada no livro, não nas apresentações
\end{frame}

%------------------------------------------------------------------------------%
\end{document}
%------------------------------------------------------------------------------%
